Při pevné vzdálenosti elektrod $d$ a ploše $S$ odpovídá zobrazený graf $J(E)$ VA-charakteristice $I(U)$: $E=U/d$, $I=JS$.
\begin{enumerate}
\item \textbf{Ohmická oblast:} slabé pole, převážná část iontů rekombinuje v plynu a jen malá část dojde k elektrodám. Koncentrace nosičů je přibližně stálá, proud roste úměrně napětí.
\item \textbf{Přechodová oblast:} s rostoucím polem více nosičů dosáhne elektrod, rekombinace ubývá. Proud již neroste přímo úměrně napětí a blíží se nasycení.
\item \textbf{Nasycení:} prakticky všechny nosiče vytvořené vnějším ionizátorem dojdou k elektrodám. Proud je určen rychlostí jejich vzniku a na napětí téměř nezávisí.
\item \textbf{Nárazová ionizace:} elektrony získají energii k ionizaci dalších molekul, vznikají elektronové laviny a proud prudce roste. Po splnění Townsendovy podmínky se výboj stane samostatným.
\end{enumerate}
Oblasti I--III a počátek IV vyžadují vnější ionizaci.
